\scanpage{38}
% archive foliation: 37

Supposons qu'on ait pu associer \uline{canoniquement}, à la catégorie de Deuring-Brauer $C$ (en termes de la seule classe d'équivalence de $C$), un objet $T$ de \unsure{$V(C)$}. Cela signifie en particulier que nous avons relevé le foncteur canonique \inserted{de Gr.-cat.}
\diagram{\begin{tikzcd}
\uline{\mathrm{Eq}}(C,T) \arrow[r] &
\uline{\mathrm{Eq}}(C) \arrow[l,bend left=30,swap,"\deleted{\unsure{a}}c"]
\end{tikzcd}}
en

Donc pour tte catégorie $C'$ équivalente à $C$, on obtient un objet $T_{\unsure{C'}}^{\unsure{c}}$ canonique dans $V(C')$.
On voit donc que pour tt $(A',C',T')$ comme plus haut, \deleted{l'ensemble} tel que $(C',T')$ est équivalent à $(C,T)_{A'}$, les équivalences en question font une \uline{catégorie-torseur} sous
\[
\uline{\mathrm{Eq}}((C,T)_{A'})
\inserted{\simeq\uline{\mathrm{Eq}}(C',T')},
\]
\deleted{donc définissent une catégorie torseur sous
$\uline{\mathrm{Eq}}(C',T')/\inserted{\operatorname{Im}}\uline{\mathrm{Eq}}(C)$ par passage}
les classes d'isom d'objets font un torseur sous
$\unsure{\operatorname{Brd}|\operatorname{Brdt}}(C',T')$,
donc définissent un torseur sous
\[
\unsure{\operatorname{Brd}|\operatorname{Brdt}}(C',T')/
\operatorname{Im}\unsure{\operatorname{Brd}|\operatorname{Brdt}}(C).
\]
Le torseur ne dépend plus du choix d'un $C$ dans la classe de catégories envisagées\ldots
Il faut donc expliciter l'application composée
\diagram{\begin{tikzcd}[column sep=large]
\unsure{\operatorname{Brd}}(C) \arrow[r,"c"] &
\unsure{\operatorname{Brd}}(C,T) \arrow[r] &
\unsure{\operatorname{Brd}}(C',T')\\
\rho \arrow[r,mapsto] &
(\rho,u_\rho) \arrow[r,mapsto] &
(\rho',u_{\rho'})\\
\unsure{P_O} \arrow[r,mapsto] \arrow[u,phantom,"\unsure{\simeq}"] &
N_{O/A}(\unsure{P_O})\unsure{\ni}u_\rho \arrow[r,mapsto] &
N_{O'/A'}(\unsure{P'_{O'}|P'_O})\unsure{\ni}(u_\rho)'
\end{tikzcd}}
